\section{Theory Questions}

\subsection{What is DMA? How does it work?}
DMA or Direct Memory Access lets peripherals communicate directly with the RAM. 
The CPU therefore does not have to copy every byte and only configures the transmission details, like source, destination, length and direction.
Then the DMA Controller performs the transmission autonomous.
The CPU is notified when the data transmission is finished or interrupted by an error and can then process the data. UART is an example for a DMA transmission, where the data is written directly into the memory buffer.


\subsection{What are the differences between Interrupt and DMA Mode?}
In interrupt-based data transfer, the peripheral signals an interrupt, prompting the CPU to handle data movement within an interrupt service routine (ISR). 
This process consumes CPU time and can become inefficient at high data rates. 
In contrast, DMA allows the CPU to configure the transfer once, after which dedicated hardware manages the data movement in the background. As a result, DMA minimizes CPU involvement and reduces the number of interrupts. These typically only triggering on buffer events or transfer completion, making it ideal for continuous or high-throughput data streams.